\scanpage{24}
on trouve donc $X\relbar\omega(X):\operatorname{Dér}(K/k_0)\relbar\mathfrak{g}_k$, application $K$-linéaire :
\[
 \omega:\operatorname{Der}(K/k_0)\longrightarrow\mathfrak{g}_K
\]
Si on veut avoir des connexions \uline{intégrables}, il faut écrire de plus la condition d'intégrabilité
\[
 \overbrace{{\Theta_0}_X\omega(Y)-{\Theta_0}_Y\omega(X)}^{\inserted{\mathrm{d}\omega(X\wedge Y)}}
 \deleted{\unsure{+[\omega(X),\omega(Y)]|-[\omega(X),\omega(Y)]}}
 =\omega([X,Y])-[\omega(X),\omega(Y)]
\]
où \deleted{\unsure{$\Theta_0$|On}} $X,Y\in\operatorname{Dér}(K/k_0)$ et ${\Theta_0}_X$, \inserted{${\Theta_0}_Y$} opèrent sur $\mathfrak{g}_K=\mathfrak{g}\otimes_{k_0}K$ via les op. sur $K$.... .

Le plus souvent, on aura en fait
\[
 \omega\in\Omega^1_{K/k_0}\otimes_K\mathfrak{g}_k
 \quad(=\Omega^1_{K/k_0}\otimes_{k_0}\mathfrak{g}_{\deleted{\unsure{K|k}}})
\]
de sorte que $\omega$ provient déjà d'un élément
\[
 \omega_{\deleted{\unsure{K}}k}\in\Omega^1_{k/k_0}\otimes_k\mathfrak{g}_k
\]
où $k$ est une ext. de type fini de $k_0$ (N.B. si $k$ est fixé, $\omega_k$ est uniquement déterminé, si on est en \uline{car. $0$}). La donnée de $\omega_k$ \unsure{revient|équivaut|révient} aussi : la donnée d'un hom $k_0$-linéaire
\[
 \Omega=\check{\mathfrak{g}}\longrightarrow\Omega^1_{k/k_0}\quad(\longrightarrow\Omega^1_{K/k_0})
\]

Nous allons examiner à quelles conditions, \deleted{\unsure{tout hom|tous les hom}} pour $V,W\in\operatorname{Ob}\mathcal{M}$ donnés, \deleted{\unsure{les}} hom \inserted{\uline{de $k_0$-vectoriels}} $u:V_{\deleted{\unsure{K}}}\longrightarrow W_{\deleted{\unsure{K}}}$ \deleted{qui} \inserted{tels que $u_K:V_K\relbar W_K$} commutent \unsure{avec|aux} $\Theta_{\unsure{x|X}}$ $(X\in\operatorname{Dér}(K/k_0))$ \deleted{\unsure{soit|sont}} commutent à $G$, en supposant $G$ connexe ($k_0$ de car. $0$). Il suffit
% archive sheet number: 23
